% §507 does not classify by how a token is *written*. It reads `cur_cmd` and
% `cur_chr`, and for a control sequence those are `eq_type` and `equiv` — its
% *meaning*. So a control sequence `\let` to a character token is that
% character, catcode and all:
%
%     @<Test if two characters match@>=
%     begin get_x_token_or_active_char;
%     if (cur_cmd>active_char)or(cur_chr>255) then {not a character}
%       begin m:=relax; n:=256;
%       end
%     else  begin m:=cur_cmd; n:=cur_chr;
%       end;
%
% `\noexpand` does not shield such a token either. §369 replaces a meaning with
% `relax`/`no_expand_flag` only when the token is expandable:
%
%     cur_cmd:=eq_type(cur_cs); cur_chr:=equiv(cur_cs);
%     if cur_cmd>max_command then
%       begin cur_cmd:=relax; cur_chr:=no_expand_flag;
%       end;
%
% hyperref depends on both halves of that. It `\futurelet`s a control sequence
% onto a space so spaces survive the `\xdef`s inside `\pdfstringdef`, then asks
% `\ifcat\relax\noexpand#1\relax` whether each token is a control sequence.
% Answering "yes" for the space sent it to `\HyPsd@CatcodeWarning` — "Token not
% allowed in a PDF string ... removing" — and every space disappeared from
% every PDF bookmark.
\input plain
\catcode`\@=11
% A control sequence \let to a character, by \let and by \futurelet.
\let\lA=A
\message{<catlet=\ifcat\lA A yes\else no\fi>}
\message{<charlet=\if\lA A yes\else no\fi>}
% Different letter, same category: \ifcat agrees, \if does not.
\message{<catAB=\ifcat\lA B yes\else no\fi>}
\message{<charAB=\if\lA B yes\else no\fi>}
% \let to a non-letter keeps that catcode.
\let\lplus=+
\message{<catplus=\ifcat\lplus +yes\else no\fi>}
\message{<catplusA=\ifcat\lplus A yes\else no\fi>}
% hyperref's idiom: \futurelet onto a space, then the control-sequence test.
% \@gobble is LaTeX's, so spell it out here.
\def\@gobble#1{}
\def\LetUS#1{\expandafter\futurelet\expandafter#1\expandafter\@gobble\space\relax}
\LetUS\usp
\message{<meaning=\meaning\usp>}
\message{<iscs=\ifcat\relax\noexpand\usp cs\else notcs\fi>}
\message{<isspace=\ifcat\usp\space yes\else no\fi>}
% \noexpand must not disturb it: the meaning is a spacer, well under
% max_command, so §369 leaves it alone.
\message{<nxspace=\ifcat\noexpand\usp\space yes\else no\fi>}
% A real \relax *is* "not a character" — cur_chr is 256.
\message{<relaxcs=\ifcat\relax\noexpand\relax cs\else notcs\fi>}
% …and so is a macro, whose meaning is above max_command and which \noexpand
% therefore does turn into relax/no_expand_flag.
\def\mac{x}
\message{<maccs=\ifcat\relax\noexpand\mac cs\else notcs\fi>}
% \chardef is *not* a character token: TeX gives it char_given, above
% active_char, so §507's first test sends it to relax/256.
\chardef\cA=`A
\message{<chardefcs=\ifcat\relax\noexpand\cA cs\else notcs\fi>}
\message{<chardefA=\ifcat\cA A yes\else no\fi>}
% Unlike almost every other scanning context, \ifcat does not skip spaces:
% get_x_token_or_active_char is a plain get_x_token. Here the space between
% `A` and `B` is the second operand, not a separator.
\message{<nospaceskip=\ifcat A B yes\else no\fi>}
\end
